\chapter{Conclusions and Suggestions for Future Research}\label{ch:conclusion}
\section{Research question revisited}
This dissertation asks whether a personalized node-private mechanism for GNN influence maximization can satisfy heterogeneous node requirements and improve Coverage over a uniform strictest-policy control. The study treats the utility claim as a testable hypothesis. It does not infer an ordering from the existence of larger budgets or from a smaller base Gaussian scale.

PDP v2 combines a fixed public identity domain and Registry with one-time personalized participation, stable overlapping slots, double clipping, and Gaussian training of all 2340 parameters. Its projection affects only the last 32 feature-map weights after each update. The original GRAT3, loss, RWR/SCS/BES sampling components, prepared graphs, and evaluation rules are retained. This makes the studied mechanism precise and independently distinguishable from historical no-projection, frozen-backbone, and rejected stability candidates.

\section{Conditional theoretical findings}
Under the stated fixed-domain, locality, ideal-randomization, and parameter-release assumptions, the mechanism has a node add/remove sensitivity bound of 12 for each slot-sum query. Full 100-step adaptive Gaussian composition and one outer amplification yield approximate target-1 bounds $(0.986940353797,2.6894142137\times10^{-6})$ and target-2/4 bounds $(1.982092738992,10^{-5})$. These satisfy the declared caps in the ideal setting.

The guarantee is conditional. It does not certify raw preprocessing, the public status of identities or preference labels, finite-precision samplers, failure and timing channels, private test graphs, or all historical research disclosures. Multiple model publication requires additional composition. A code test or an arithmetic certificate cannot remove those assumptions.

The target-4 class has the same participation and bound as target 2. The mechanism therefore uses some heterogeneity but not all available allowance. Future claims about improved budget utilization must identify a new valid mechanism rather than interpreting the unused allowance as already realized utility.

\section{Independent empirical findings}
The final study contains 2900 independent-confirmation analysis units across the fixed Email, Bitcoin, LastFM, HepPh, and Facebook graphs. Per graph, it includes 200 PDP v2, 200 strict Uniform1, 140 relaxed Uniform4, and 40 Non-Private outcomes. The complete execution record retains 2902 fit attempts, two fit failures, authorized recovery lineage, and all low-utility results. Attempts are not additional statistical units.

The equal-weight PDP--Uniform1 estimate is $+0.7845395873$ pp. Its one-sided $p$-value is $0.1874667289$, the one-sided 95\% lower bound is $-0.6702897323$ pp, and the two-sided 95\% interval is $[-0.9492939081,2.518373083]$ pp. The study does not establish the proposed positive average effect at the frozen threshold, or an improvement of the planned 3 pp magnitude. It also does not establish equivalence or exclude every potentially useful smaller benefit.

There are 103 severe-low-utility outcomes among 1000 PDP units, compared with 115 among 1000 strict Uniform units. Facebook has more observed PDP low-tail events and a larger sample SD than the strict control. The combined evidence does not establish a general stability or risk advantage. Utility gaps to Uniform4 and Non-Private remain, but their different protection roles prevent treating them as compliant substitutes for all heterogeneous policies.

\section{Implications for private graph learning}
The main practical lesson is that a feasible privacy construction and lower nominal noise need not yield reliable utility improvement. Information removed by participation, bounded slot contributions, high-dimensional optimization, and nonlinear seed ranking may interact. The available record identifies this gap without resolving all of its causes.

The historical frozen-backbone results and Gaussian-path diagnostics strengthen the motivation to study training-space geometry and random-path sensitivity. They do not prove a unique explanation. Retaining unsuccessful G3/C3 candidates and adverse graph results prevents a misleading narrative in which every development change is retrospectively declared an improvement.

A contribution can therefore be empirical characterization rather than performance victory. The full independent matrix provides a bounded answer to a concrete question, quantifies uncertainty, and identifies limitations that future mechanisms must address. This is stronger evidence than presenting a favorable initialization subset or treating failure to detect significance as proof of no effect.

\section{Future research directions}
One direction is better use of heterogeneous budgets. A mechanism that distinguishes targets 2 and 4 might permit different contributions or accounting paths, but changing probabilities, noise, or sensitivities requires a new proof. Lowering the common Gaussian scale to the relaxed reference value would not preserve the existing target-2 and strict-node guarantees automatically.

A second direction is optimization stability. Reduced trainable subspaces, improved public representations, or carefully analyzed update constraints may help, but the current frozen-backbone evidence is developmental. Any proposal should be evaluated against fixed graph-specific utility and low-tail criteria, with all trajectories retained. Previously rejected averaging or row-sum candidates should not be restarted merely because their aggregate averages looked promising.

A third direction is evaluation breadth. New graph families, repeated prepared splits, alternative policy proportions, and stochastic multi-step diffusion could test whether the current conclusion persists. Such work would broaden the target of inference; it would not retroactively make the present five-graph estimate representative of all networks.

A fourth direction is complete release analysis. A deployable pipeline needs justified public-domain construction, protection for changing target graphs and selected identities, composition for model collections, and a treatment of numerical and execution channels. These are distinct tasks from improving Coverage. Future research should make their relationship explicit rather than treating a training accountant as a system-wide certificate.

The conditional privacy construction is supported within its assumptions. Independent evaluation does not establish a Coverage or stability advantage for PDP v2 over the strict uniform control. Privacy feasibility and empirical utility improvement therefore remain distinct findings.
